%==============================================================================
% tcc.tex --- Avaliacao do custo energetico de politicas de offloading
% baseadas em confianca em cascatas hierarquicas de inferencia de LLMs
% sobre redes opticas passivas
%
% Rascunho parcial: capitulos 1 (Introducao) e 2 (Revisao Bibliografica).
% Metodologia, resultados e conclusao ainda serao escritos.
%
% Compilar com:
%   pdflatex tcc
%   bibtex   tcc
%   pdflatex tcc
%   pdflatex tcc
%==============================================================================

\documentclass[cic,dipl]{infufrgs}

\usepackage[utf8]{inputenc}
\usepackage{graphicx}
\usepackage{amsmath}

%------------------------------------------------------------------------------
% Metadados do documento
%------------------------------------------------------------------------------
\title{Avalia\c{c}\~ao do custo energ\'etico de pol\'iticas de offloading
  baseadas em confian\c{c}a em cascatas hier\'arquicas de infer\^encia de
  LLMs sobre redes \'opticas passivas}

\author{Amorim}{Diego Hommerding}

% Orientador ainda pendente de confirmacao formal (ver documento mestre,
% secao 10 e 11, item 1). Nome completo e titulacao a confirmar.
\advisor[Prof.]{Nazar}{}

%------------------------------------------------------------------------------
% Palavras-chave
%------------------------------------------------------------------------------
\keyword{Infer\^encia distribu\'ida de LLMs}
\keyword{Redes \'opticas passivas}
\keyword{Efici\^encia energ\'etica}
\keyword{Offloading baseado em confian\c{c}a}
\keyword{Computa\c{c}\~ao em n\'evoa}

%==============================================================================
\begin{document}
%==============================================================================

\maketitle

%------------------------------------------------------------------------------
\begin{abstract}
Arquiteturas hier\'arquicas de infer\^encia de LLMs distribuem o
processamento entre dispositivo de usu\'ario, borda, fog e nuvem, decidindo
em cada n\'ivel se respondem localmente ou encaminham a requisi\c{c}\~ao
para um modelo maior. As pol\'iticas atuais dessa decis\~ao, como o
RecServe \cite{wu2025recserve} e a arquitetura sobre TWDM-PON de Pakpahan e
Hwang \cite{pakpahan2026continuum}, otimizam qualidade, lat\^encia e uso de
banda, mas n\~ao consideram o consumo de energia. Este trabalho avalia o
impacto energ\'etico dessas pol\'iticas de offloading, estabelecendo uma
base de medi\c{c}\~ao compar\'avel entre as camadas e caracterizando o
custo de cada uma a partir de dados j\'a publicados. Com isso, busca
identificar quando uma decis\~ao de roteamento ciente de energia
diferiria da decis\~ao baseada apenas em confian\c{c}a, apontando
caminhos para uma infer\^encia distribu\'ida mais eficiente sobre redes
de acesso \'optico.
\end{abstract}

\tableofcontents

\listoffigures

\listoftables

\begin{listofabbrv}{TWDM-PON}
  \item[LLM]  Large Language Model
  \item[NPU]  Neural Processing Unit
  \item[ONU]  Optical Network Unit
  \item[OLT]  Optical Line Terminal
  \item[PON]  Passive Optical Network
  \item[PUE]  Power Usage Effectiveness
  \item[TWDM-PON] Time and Wavelength Division Multiplexed Passive Optical Network
\end{listofabbrv}

%==============================================================================
\chapter{Introdu\c{c}\~ao}
\label{cap:introducao}
%==============================================================================

\section{Contextualiza\c{c}\~ao}

Arquiteturas hier\'arquicas de infer\^encia de modelos de linguagem de
grande porte (LLMs) distribuem o processamento entre m\'ultiplos n\'iveis:
o dispositivo do usu\'ario, a borda da rede, a camada de fog e a nuvem. Em
cada n\'ivel, o sistema decide se responde localmente, usando um modelo
menor e mais barato, ou encaminha a requisi\c{c}\~ao para um modelo maior
em um n\'ivel superior, mais caro por\'em mais capaz. Essa decis\~ao de
roteamento, chamada de offloading, \'e o n\'ucleo de dois trabalhos
recentes que servem de baseline para esta pesquisa.

O RecServe \cite{wu2025recserve} prop\~oe um esquema de offloading
recursivo baseado em um limiar de confian\c{c}a: em cada camada, a
confian\c{c}a da resposta gerada localmente \'e comparada contra um
limiar din\^amico, calculado como um quantil de um hist\'orico local de
confian\c{c}a mantido em uma janela deslizante. A requisi\c{c}\~ao \'e
escalonada para a camada seguinte somente quando essa confian\c{c}a fica
abaixo do limiar. O limiar \'e din\^amico de prop\'osito, para n\~ao
depender de um corte fixo escolhido manualmente, que n\~ao generalizaria
conforme a distribui\c{c}\~ao de confian\c{c}a do modelo muda com o
tr\'afego.

Pakpahan e Hwang \cite{pakpahan2026continuum} prop\~oem uma arquitetura de
infer\^encia hier\'arquica sobre redes \'opticas passivas com
multiplexa\c{c}\~ao por divis\~ao de tempo e comprimento de onda
(TWDM-PON), mapeando cada camada da hierarquia, usu\'ario, ONU (Optical
Network Unit), OLT (Optical Line Terminal) e nuvem, a uma faixa de
tamanho de modelo e a uma classe de hardware, de forma a atender
requisitos de qualidade de servi\c{c}o sobre a rede de acesso \'optico.

\section{Problema e lacuna}

Ambos os trabalhos otimizam aspectos relevantes do sistema, mas nenhum
considera energia como crit\'erio de decis\~ao. O RecServe otimiza a
quantidade de bytes trocados durante a cascata de offloading. Pakpahan e
Hwang otimizam qualidade de servi\c{c}o sobre a rede \'optica,
especificamente lat\^encia, jitter e vaz\~ao. O consumo de energia por
decis\~ao de roteamento n\~ao aparece como vari\'avel de otimiza\c{c}\~ao
em nenhum dos dois.

Essa lacuna \'e relevante porque a intui\c{c}\~ao mais comum, a de que
responder localmente \'e sempre mais eficiente energeticamente do que
escalar para um modelo maior, n\~ao \'e necessariamente verdadeira quando
se considera a efici\^encia energ\'etica do hardware de cada camada e o
regime de opera\c{c}\~ao (batelada ou n\~ao) em que ele processa
requisi\c{c}\~oes. Um hardware de nuvem, mesmo servindo um modelo muito
maior, pode ser mais eficiente por token do que um hardware de fog mais
fraco servindo um modelo intermedi\'ario, dependendo de como ambos s\~ao
utilizados em produ\c{c}\~ao.

A pergunta central deste trabalho \'e: considerar energia, e n\~ao s\'o
confian\c{c}a, desloca o ponto de decis\~ao de roteamento em uma cascata
hier\'arquica de infer\^encia de LLMs? Essa pergunta se sustenta mesmo
sem uma pol\'itica de roteamento nova: medir energia sobre a arquitetura
de offloading j\'a \'e uma contribui\c{c}\~ao, na medida em que nenhum
dos baselines mede ou otimiza esse recurso. Uma pol\'itica de decis\~ao
ciente de energia, a ser detalhada em cap\'itulos futuros deste
documento, \'e um resultado adicional, n\~ao uma depend\^encia da
contribui\c{c}\~ao principal.

\section{Objetivos}

O objetivo geral deste trabalho \'e avaliar o impacto energ\'etico de
pol\'iticas de offloading baseadas em confian\c{c}a em cascatas
hier\'arquicas de infer\^encia de LLMs sobre redes \'opticas passivas.

Os objetivos espec\'ificos s\~ao:

\begin{enumerate}
  \item Estabelecer uma base de medi\c{c}\~ao de energia compar\'avel
    entre as camadas de usu\'ario, borda, fog e nuvem, a partir de dados
    j\'a publicados na literatura, fixando os fatores que impedem a
    compara\c{c}\~ao direta entre estudos (batch, precis\~ao, fase da
    infer\^encia, fronteira de medi\c{c}\~ao e efici\^encia do
    datacenter).
  \item Caracterizar o custo energ\'etico de cada camada da hierarquia
    usando essa base, identificando onde a literatura j\'a oferece
    medi\c{c}\~ao direta no hardware da classe correspondente e onde
    permanecem lacunas.
  \item Instrumentar o RecServe para contabilizar energia por decis\~ao
    de roteamento, associando cada trecho da cascata percorrida por uma
    requisi\c{c}\~ao ao custo energ\'etico da camada que a processou.
  \item Propor e avaliar uma pol\'itica de roteamento ciente de energia,
    comparando-a com a pol\'itica original baseada apenas em confian\c{c}a,
    e caracterizar em quais condi\c{c}\~oes as duas decis\~oes divergem.
\end{enumerate}

\section{Organiza\c{c}\~ao do texto}

O restante deste documento est\'a organizado da seguinte forma. O
Cap\'itulo~\ref{cap:revisao} apresenta a revis\~ao bibliogr\'afica,
detalhando os dois trabalhos de baseline, as m\'etricas de energia
usadas na literatura de infer\^encia de LLMs e os estudos de medi\c{c}\~ao
de energia dispon\'iveis para cada classe de hardware da hierarquia. Os
cap\'itulos de metodologia, resultados e conclus\~ao ainda ser\~ao
escritos, \`a medida que a base de medi\c{c}\~ao for fechada e o
experimento de compara\c{c}\~ao entre pol\'iticas for executado.

%==============================================================================
\chapter{Revis\~ao Bibliogr\'afica}
\label{cap:revisao}
%==============================================================================

Este cap\'itulo apresenta os trabalhos que fundamentam esta pesquisa. A
Se\c{c}\~ao~\ref{sec:recserve} detalha o RecServe, o baseline de
implementa\c{c}\~ao da pol\'itica de offloading. A
Se\c{c}\~ao~\ref{sec:pakpahan} detalha a arquitetura de infer\^encia
hier\'arquica sobre TWDM-PON de Pakpahan e Hwang, usada como refer\^encia
de mapeamento entre camada, modelo e hardware. A
Se\c{c}\~ao~\ref{sec:metricas} apresenta as m\'etricas de energia
propostas na literatura para infer\^encia de LLMs. A
Se\c{c}\~ao~\ref{sec:medicoes} revisa os estudos de medi\c{c}\~ao de
energia dispon\'iveis para cada classe de hardware da hierarquia. A
Se\c{c}\~ao~\ref{sec:sintese} sintetiza a lacuna identificada nesses
trabalhos, que sustenta a contribui\c{c}\~ao desta pesquisa.

\section{Offloading recursivo baseado em confian\c{c}a: RecServe}
\label{sec:recserve}

O RecServe \cite{wu2025recserve} \'e um esquema de offloading para
servi\c{c}o de LLMs em redes de m\'ultiplas camadas. Para uma consulta
$x$ processada na camada $M$ com tarefa $\tau$, a decis\~ao segue a
regra

\begin{equation}
  \text{Decis\~ao}(x) =
  \begin{cases}
    M, & \text{se } C_{M,\tau}(x) \geq T_{M,\tau}(\beta) \\
    \text{escalona para } M+1, & \text{se } C_{M,\tau}(x) < T_{M,\tau}(\beta)
  \end{cases}
\end{equation}

em que $C_{M,\tau}(x)$ \'e a confian\c{c}a da resposta gerada na camada
$M$ e $T_{M,\tau}(\beta)$ \'e o quantil $\beta$ sobre o hist\'orico local
de confian\c{c}a $H_{M,\tau}$, mantido em uma fila deslizante de tamanho
$k$. A cascata escalona apenas um degrau por vez, nunca pulando
camadas, e a decis\~ao \'e tomada localmente em cada camada, sem
depend\^encia de informa\c{c}\~ao das demais.

O limiar din\^amico \'e a caracter\'istica central do RecServe: em vez de
um corte fixo escolhido manualmente, o sistema recalcula o limiar a
partir da distribui\c{c}\~ao recente de confian\c{c}a observada na
pr\'opria camada, adaptando-se a mudan\c{c}as no tr\'afego ou na
dificuldade das consultas ao longo do tempo. Outras partes do sistema,
no entanto, permanecem est\'aticas por escolha de projeto: o modelo
usado em cada camada \'e uma configura\c{c}\~ao fixa, e os
hiperpar\^ametros $\beta$ e $k$ s\~ao fixados por experimento, n\~ao
reestimados durante a execu\c{c}\~ao. O c\'odigo do RecServe \'e aberto,
o que permite instrument\'a-lo diretamente para os fins desta pesquisa.

O RecServe otimiza a quantidade de bytes de comunica\c{c}\~ao trocados
durante a cascata de offloading. Energia n\~ao \'e uma vari\'avel
considerada na decis\~ao de roteamento nem contabilizada como m\'etrica
de avalia\c{c}\~ao.

\section{Infer\^encia hier\'arquica sobre TWDM-PON}
\label{sec:pakpahan}

Pakpahan e Hwang \cite{pakpahan2026continuum} prop\~oem uma arquitetura
de infer\^encia hier\'arquica definida por software sobre uma rede
\'optica passiva com multiplexa\c{c}\~ao por divis\~ao de tempo e
comprimento de onda (TWDM-PON). A proposta mapeia cada n\'ivel da
hierarquia de rede a uma faixa de tamanho de modelo e a uma classe de
hardware, conforme resumido na Tabela~\ref{tab:pakpahan}.

\begin{table}[htbp]
  \centering
  \caption{Mapeamento entre camada, modelo e hardware proposto por Pakpahan e Hwang \cite{pakpahan2026continuum}.}
  \label{tab:pakpahan}
  \begin{tabular}{lll}
    \hline
    \textbf{Camada} & \textbf{Modelo} & \textbf{Hardware} \\
    \hline
    Usu\'ario & tiny-LLM (menos de 1B) & NPU de celular \\
    ONU (borda) & 0,5 a 7B & Jetson \\
    OLT (fog) & 7 a 13B & A30 / T4 \\
    Nuvem & 70B ou mais & A100 / H100 \\
    \hline
  \end{tabular}
\end{table}

Essa arquitetura otimiza qualidade de servi\c{c}o sobre a rede \'optica,
particularmente lat\^encia, jitter e vaz\~ao entre as camadas. \'E usada
neste trabalho como arquitetura de refer\^encia para o mapeamento entre
camada e classe de hardware, e como ponto de partida para o crit\'erio de
sele\c{c}\~ao de fontes de medi\c{c}\~ao de energia apresentado na
Se\c{c}\~ao~\ref{sec:medicoes}. Assim como o RecServe, essa arquitetura
n\~ao considera consumo de energia entre os crit\'erios de otimiza\c{c}\~ao.

\section{M\'etricas de energia para infer\^encia de LLMs}
\label{sec:metricas}

Dois trabalhos recentes propõem m\'etricas de energia especificamente
para infer\^encia de LLMs e s\~ao usados como refer\^encia de m\'etodo
nesta pesquisa. Solovyeva e Castor \cite{solovyeva2026green} prop\~oem
medir o custo energ\'etico de infer\^encia em joules por token de sa\'ida
(J/token), separando a medi\c{c}\~ao por fase de infer\^encia (prefill e
decode). Essa \'e a unidade adotada como base para a compara\c{c}\~ao
entre camadas neste trabalho, descrita em detalhe no cap\'itulo de
metodologia.

Li e Lu \cite{li2026ecothink} prop\~oem uma m\'etrica de emiss\~ao de
carbono por consulta (gCO2/query), derivada da energia por consulta
multiplicada pela efici\^encia do datacenter (\emph{Power Usage
Effectiveness}, PUE) e pela intensidade de carbono da rede el\'etrica
local. Essa m\'etrica \'e mais pr\'oxima do objetivo final de avaliar
impacto ambiental, mas depende de fatores externos ao hardware
(localiza\c{c}\~ao do datacenter, matriz energ\'etica local) que
dificultam a compara\c{c}\~ao direta entre camadas heterog\^eneas de uma
cascata distribu\'ida. Este trabalho usa J/token como unidade prim\'aria
e deriva gCO2/query como m\'etrica secund\'aria, quando aplic\'avel.

\section{Medi\c{c}\~ao de energia por classe de hardware}
\label{sec:medicoes}

Esta se\c{c}\~ao revisa os estudos de medi\c{c}\~ao de energia
dispon\'iveis para cada classe de hardware da hierarquia definida por
Pakpahan e Hwang \cite{pakpahan2026continuum}. Nem todos os estudos usam
exatamente o mesmo hardware citado por Pakpahan e Hwang; quando a
medi\c{c}\~ao exata n\~ao est\'a dispon\'ivel, prioriza-se o estudo mais
rigoroso dentro da mesma classe de implanta\c{c}\~ao (dispositivo de
usu\'ario, borda, fog ou hyperscale), crit\'erio detalhado no cap\'itulo
de metodologia.

\subsection{Dispositivo de usu\'ario}

Cai et al. \cite{cai2026npu} medem energia por token e vaz\~ao em um
processador Snapdragon~8~Elite~Gen~5, separando a medi\c{c}\~ao por fase
de infer\^encia e por unidade de computa\c{c}\~ao (CPU, GPU e NPU), com
quantiza\c{c}\~ao de 4 bits e tamanho de lote igual a 1. Um achado
relevante desse estudo \'e que, no decode, a CPU do celular costuma ser
mais eficiente energeticamente do que a GPU, e a NPU pode consumir de 2 a
4 vezes mais energia de decode apesar de ser frequentemente associada a
efici\^encia energ\'etica; a NPU domina apenas na fase de prefill.

\subsection{Borda}

Kubwimana e Huang \cite{kubwimana2025edgereasoning} medem energia por
token em uma placa NVIDIA Jetson AGX Orin, o hardware exato atribu\'ido
\`a ONU por Pakpahan e Hwang, separando prefill e decode, em FP16 e
quantiza\c{c}\~ao de 4 bits, com tamanho de lote igual a 1. O estudo
confirma que o decode domina 99\% ou mais do tempo e da energia total de
infer\^encia nessa classe de hardware, consistente com um regime
limitado por banda de mem\'oria (\emph{memory-bound}) para infer\^encia
com lote pequeno.

\subsection{Fog}

Tr\^es estudos recentes medem energia em GPUs de classe fog (A30, T4,
L4, L40S). Fadel Argerich, F\"urst e Pati\~no-Mart\'inez
\cite{fadelargerich2026wattcounts} medem energia por token para
diversos modelos e GPUs em cen\'ario de servidor com lote alto,
incluindo um ponto direto em uma placa A30. Em trabalho complementar,
os mesmos autores \cite{fadelargerich2026wattgpu} medem pot\^encia sob
restri\c{c}\~ao de lat\^encia entre tokens (\emph{inter-token latency}),
tamb\'em em GPUs A30. Wilhelm, Wittkopp e Kao
\cite{wilhelm2026energypertoken} defendem o uso de energia por token
como m\'etrica prim\'aria de avalia\c{c}\~ao de infer\^encia e
apresentam medi\c{c}\~oes em GPU L40S com tamanho de lote igual a 1,
o cen\'ario mais pr\'oximo da base de medi\c{c}\~ao adotada nas camadas
de usu\'ario e borda deste trabalho.

\subsection{Nuvem}

Samsi et al. \cite{samsi2023words2watts} medem energia por token para
modelos da fam\'ilia Llama de 65 a 70 bilh\~oes de par\^ametros em GPUs
A100, com tamanho de lote igual a 64. Oviedo et al.
\cite{oviedo2026joule}, publicado no peri\'odico \emph{Joule}, comparam
energia por consulta entre diferentes sistemas de \emph{serving} em
produ\c{c}\~ao para o modelo Llama~3.1~70B, incluindo cen\'arios com e
sem \emph{batching}, mostrando uma varia\c{c}\~ao de at\'e uma ordem de
grandeza entre o regime de produ\c{c}\~ao real e um cen\'ario isolado
sem otimiza\c{c}\~oes de \emph{serving}.

Um achado consistente entre os estudos revisados nesta subse\c{c}\~ao,
tamb\'em observado nos estudos de fog citados acima, \'e que nenhum
registra medi\c{c}\~ao de fog ou nuvem em tamanho de lote igual a 1 que
reflita um cen\'ario de produ\c{c}\~ao real: essas camadas, na pr\'atica,
sempre atendem m\'ultiplas requisi\c{c}\~oes simultaneamente para
amortizar o custo do hardware. Essa aus\^encia \'e tratada neste
trabalho como uma caracter\'istica estrutural do dom\'inio, discutida em
detalhe no cap\'itulo de metodologia.

\section{S\'intese}
\label{sec:sintese}

Os dois baselines revisados neste cap\'itulo, RecServe
\cite{wu2025recserve} e a arquitetura sobre TWDM-PON de Pakpahan e Hwang
\cite{pakpahan2026continuum}, otimizam qualidade, lat\^encia e uso de
banda, mas nenhum considera energia como crit\'erio de decis\~ao. As
m\'etricas de energia propostas na literatura \cite{solovyeva2026green,
li2026ecothink} oferecem o vocabul\'ario metodol\'ogico para medir esse
custo, e os estudos de medi\c{c}\~ao por classe de hardware
\cite{cai2026npu, kubwimana2025edgereasoning,
fadelargerich2026wattcounts, fadelargerich2026wattgpu,
wilhelm2026energypertoken, samsi2023words2watts, oviedo2026joule}
oferecem dados publicados suficientes para caracterizar o custo
energ\'etico de cada camada da hierarquia, ainda que com graus
desiguais de cobertura entre as camadas. \'E essa combina\c{c}\~ao,
crit\'erios de decis\~ao sem energia e dados de energia dispersos na
literatura sem uma base comparável entre camadas, que sustenta a
contribui\c{c}\~ao desta pesquisa, detalhada no cap\'itulo de
metodologia.

%==============================================================================
% Referencias
%==============================================================================
\bibliographystyle{abntex2-alf}
\bibliography{tcc}

\end{document}
